\documentclass[../main]{subfiles}
\begin{document}
\newpage

\subsubsection{\texttt{armax\_filtered}}

\begin{lstlisting}
theta, m = pysib.armax_filtered(u, y, na, nb, nc, nz)
\end{lstlisting}

Identifies an ARMAX model using a cost-function shaping strategy to
improve convergence.  The model structure and the criterion minimized
at each stage are the same as for \texttt{armax}:

\[
  A(q^{-1})\,y(t) = B(q^{-1})\,u(t-n_k) + C(q^{-1})\,e(t), \qquad
  V(\theta) = \frac{1}{N}\sum_t \varepsilon(t;\theta)^2.
\]

Before the final unfiltered ARMAX step, 9 auxiliary ARMAX problems
are solved on filtered versions of the data.  At stage $i$, both
\texttt{u} and \texttt{y} are passed through the same first-order
low-pass filter sequence used by \texttt{oe\_filtered}:
\[
  a_i = 0.05^{1/\tau_i}, \qquad
  \tau_i \in \{36,32,28,24,20,16,12,8,4\}.
\]
The ARMAX criterion is
minimized on the filtered data, progressively widening the passband
and reshaping the cost function to reduce the attraction of
undesirable local minima before the estimate is refined on the
original data.

\subsubsection*{Arguments}
\begin{center}
\begin{tabular}{lll}
  \toprule
  Name & Type & Description \\
  \midrule
  \texttt{u}  & array\_like & Input signal. \\
  \texttt{y}  & array\_like & Output signal. \\
  \texttt{na} & int & Number of $A$ parameters. \\
  \texttt{nb} & int & Number of $B$ parameters. \\
  \texttt{nc} & int & Number of $C$ parameters. \\
  \texttt{nz} & int & Input delay $n_k \geq 0$. \\
  \bottomrule
\end{tabular}
\end{center}

\subsubsection*{Returns}
\begin{center}
\begin{tabular}{lll}
  \toprule
  Name & Type & Value \\
  \midrule
  \texttt{theta} & ndarray & $[a_1,\ldots,a_{n_a},\; b_0,\ldots,b_{n_b-1},\; c_1,\ldots,c_{n_c}]$ \\
  \texttt{m}     & dict    & \texttt{A}, \texttt{B}, \texttt{C} free; \texttt{D}$=$\texttt{F}$=[1]$ \\
  \bottomrule
\end{tabular}
\end{center}

\end{document}
